% Owner: Kais
\section{Beacon Deposition Decision}

Beacons do two jobs: they mark where an event is, and they keep the chain of
radio links back to the PEB unbroken, so that every event reaches the Command
Post and an Executor can follow the chain. The Writer therefore drops a
beacon in two cases.

\paragraph{1. An event is detected.} The Writer stops at the event and drops
an \emph{event beacon} there (gas, fire, victim alive or dead), with the
severity and the path from its previous beacon. The beacon sits at the spot
the Executor must reach, so the last primitive of the path leads it exactly
there.

\paragraph{2. The Writer is about to lose the mesh.} The Writer runs an
ESP32, whose ESP-NOW receive callback reports the signal strength (RSSI) of
every frame. While exploring, it keeps hearing the 3\,s rebroadcasts of the
beacons behind it. When the strongest of them falls below a threshold
(about $-80$\,dBm \todo{calibrate}), it drops a \emph{connectivity beacon}
(event type \texttt{0x00}). The threshold sits above the level where ESP-NOW
stops working (around $-90$\,dBm), so the new beacon still has a reliable
link to the previous one. This uses beacons only where the building actually
weakens the signal (walls, corners, long corridors), instead of at a fixed
spacing that would waste them in open areas.

\paragraph{Loading and dropping.} Before a drop, the Writer sends the frame
by ESP-NOW unicast to the next beacon in its magazine. The beacon stores it
in non-volatile memory, turns its LED on and starts broadcasting; the
mechanism then releases it (Section~\ref{sec:mechanics}). The Writer moves on
without waiting for a confirmation; a beacon that fails is detected later by
the Command Post (Section~\ref{sec:failure-cases}).

\paragraph{Beacon budget.} A Writer carries $K$ beacons \todo{K, set by the
magazine}. When only 2 remain, it keeps them for events, stops exploring new
branches, and returns to the PEB to reload.
